\chapter{線形代数II：固有値・内積・標準形}
\label{ch:tb-linear-two}

\section{固有値と対角化}

\begin{definition}{固有値・固有空間}{tb-eigenvalue}
$Av=\lambda v$ を満たす $v\ne0$ が存在するとき、$\lambda$ を $A$ の固有値、
$v$ を固有ベクトルという。
\[
E_\lambda=\ker(A-\lambda I)
\]
を固有空間という。固有多項式を
$\chi_A(t)=\det(tI-A)$ とする。
\end{definition}

固有値と固有空間は
\begin{enumerate}
  \item $\chi_A(t)$ を因数分解して固有値を得る。
  \item 各 $\lambda$ で $(A-\lambda I)v=0$ を行基本変形し、基底を求める。
\end{enumerate}
の順で計算する。$\det A$ は固有値の積、$\tr A$ は固有値の和である
（重複度を込める）。

\begin{definition}{対角化}{tb-diagonalization}
正則行列 $P=(v_1\ \cdots\ v_n)$ に対し
\[
P^{-1}AP=\diag(\lambda_1,\ldots,\lambda_n)
\]
となることを対角化という。これは $v_1,\ldots,v_n$ が固有ベクトルからなる基底であることと同値である。
\end{definition}

固有値 $\lambda$ の固有多項式での重複度を代数的重複度 $m_a(\lambda)$、
$\dim E_\lambda$ を幾何的重複度 $m_g(\lambda)$ といい
\[
1\le m_g(\lambda)\le m_a(\lambda).
\]
固有多項式が基礎体上で一次式に分解するとき、
全ての固有値で両者が等しいことと対角化可能であることは同値である。
相異なる固有値の固有ベクトルは一次独立なので、固有値が $n$ 個あれば自動的に対角化できる。

\begin{proposition}{行列が満たす多項式の使い方}{tb-matrix-polynomial}
$p(A)=O$ なら、$A$ の固有値 $\lambda$ は全て $p(\lambda)=0$ を満たす。
さらに
\[
p(t)=\prod_{i=1}^r(t-\lambda_i)\qquad(\lambda_i\ne\lambda_j\ (i\ne j))
\]
と相異なる一次式に分解できれば、$A$ は対角化可能である。このとき固有値 $\lambda_i$ の成分は
\[
P_i x=\prod_{j\ne i}\frac{A-\lambda_jI}{\lambda_i-\lambda_j}\,x
\]
で取り出せ、$x=\sum_iP_ix$、$P_ix\in E_{\lambda_i}$ となる。
\end{proposition}

$Av=\lambda v$ なら $p(A)v=p(\lambda)v$ なので前半が従う。
後半は Lagrange 補間の恒等式 $\sum_iP_i=I$ と
$(A-\lambda_iI)P_i=O$ を用いればよい。

対角化できれば
\[
A^k=PD^kP^{-1},\qquad
D^k=\diag(\lambda_1^k,\ldots,\lambda_n^k).
\]
行列で表された漸化式 $X_n=AX_{n-1}$ も $X_n=A^{n-n_0}X_{n_0}$ として同じ方法で解く。

\section{冪零行列と Jordan 標準形}

\begin{definition}{冪零と Jordan ブロック}{tb-jordan}
$N^m=O$ となる正整数 $m$ が存在するとき $N$ を冪零という。
\[
J_k(\lambda)=
\begin{pmatrix}
\lambda&1&&0\\
&\lambda&\ddots&\\
&&\ddots&1\\
0&&&\lambda
\end{pmatrix}
\]
を Jordan ブロックといい、これらを対角に並べた行列を Jordan 標準形という。
\end{definition}

冪零行列の固有値は0だけである。実際 $Nv=\lambda v$ なら
$0=N^mv=\lambda^mv$ なので $\lambda=0$ である。
したがって全ての $k\ge1$ で $\tr N^k=0$ となる。

\begin{proposition}{Jordan ブロックの読み方}{tb-jordan-nullities}
固有値 $\lambda$ について $N=A-\lambda I$ とする。
\begin{itemize}
  \item ブロックの個数は $\dim\ker N$。
  \item 大きさ $j$ 以上のブロックの個数は
        $\dim\ker N^j-\dim\ker N^{j-1}$。
  \item 最大ブロックの大きさは $N^j$ が一般化固有空間上で0になる最小の $j$。
\end{itemize}
\end{proposition}

3次で固有値が一つなら、$\dim\ker N=3,2,1$ に応じてブロック型は
$1+1+1$, $2+1$, または $3$ である。
パラメータを含む問題では、固有値が衝突する値を先に分けて核の次元を調べる。

\section{不変部分空間と上三角化}

\begin{definition}{不変部分空間}{tb-invariant-subspace}
$W\subset V$ が $A$ 不変であるとは $AW\subset W$ をいう。
1次元不変部分空間は固有ベクトル $v$ が張る直線 $\operatorname{span}\{v\}$ に限る。
\end{definition}

$V=\R^3$ の2次元部分空間は $W=\ker \ell$（$\ell\ne0$ は行ベクトル）と書ける。
このとき
\[
AW\subset W
\Longleftrightarrow \ell A=\lambda\ell\text{ となる }\lambda\text{ が存在する}.
\]
したがって2次元不変部分空間は $A^{\mathsf T}$ の固有ベクトル $u$ を求め、
$W=\{x\mid u^{\mathsf T}x=0\}$ とすれば分類できる。

\begin{theorem}{不変旗と上三角化}{tb-invariant-flag}
基底 $v_1,\ldots,v_n$ と
$W_j=\operatorname{span}\{v_1,\ldots,v_j\}$ が全て $AW_j\subset W_j$ を満たすとする。
$P=(v_1\ \cdots\ v_n)$ とおけば $P^{-1}AP$ は上三角行列である。
\end{theorem}

\begin{proof}
$Av_j\in W_j$ なので
$Av_j=\sum_{i=1}^j t_{ij}v_i$ と書ける。
これは $AP=PT$ の第 $j$ 列であり、$i>j$ の成分 $t_{ij}$ が0、すなわち $T$ が上三角であることを意味する。
\end{proof}

\section{内積、直交行列、ユニタリ行列}

実内積を $\inner{x}{y}=x^{\mathsf T}y$、複素内積を $\inner{x}{y}=x^*y$、
$\norm{x}=\sqrt{\inner{x}{x}}$ とする。
実内積はノルムから
\[
\inner{x}{y}=\frac12\left(\norm{x+y}^2-\norm{x}^2-\norm{y}^2\right)
\]
と復元できる。

一次独立な $v_1,\ldots,v_m$ は Gram--Schmidt 法
\[
u_1=v_1,\qquad
u_k=v_k-\sum_{j=1}^{k-1}
\frac{\inner{u_j}{v_k}}{\inner{u_j}{u_j}}u_j,\qquad
e_k=\frac{u_k}{\norm{u_k}}
\]
で正規直交系 $e_1,\ldots,e_m$ にできる。
固有空間の中でこの操作を行っても、得られるベクトルは同じ固有空間に残る。

\begin{definition}{直交行列・ユニタリ行列}{tb-orthogonal-unitary}
実行列 $Q$ が $Q^{\mathsf T}Q=I$ を満たすとき直交行列、
複素行列 $U$ が $U^*U=I$ を満たすときユニタリ行列という。
列ベクトルが正規直交基底であることと同値であり
\[
Q^{-1}=Q^{\mathsf T},\qquad U^{-1}=U^*.
\]
\end{definition}

\begin{proposition}{長さ・内積・直交性の同値}{tb-isometry-equivalence}
実正方行列 $Q$ について次は同値である。
\[
Q^{\mathsf T}Q=I,\qquad
\norm{Qx}=\norm x\ (\forall x),\qquad
\inner{Qx}{Qy}=\inner{x}{y}\ (\forall x,y).
\]
\end{proposition}

\begin{proof}
最初から最後は $\inner{Qx}{Qy}=x^{\mathsf T}Q^{\mathsf T}Qy$ で従う。
長さの保存から内積の保存は上の偏極公式を使う。
内積の保存からは
$x^{\mathsf T}(Q^{\mathsf T}Q-I)y=0$ が全ての $x,y$ で成り立つため $Q^{\mathsf T}Q=I$ である。
\end{proof}

3次直交行列の二列が分かっているとき、残る単位列ベクトルはその外積の $\pm$ である。
最後に行列式の符号や与えられた成分でどちらかを選ぶ。
また $(Q^{\mathsf T})^k=Q^{-k}$ なので、直交行列の複雑な積は指数法則へ直す。

\section{対称・交代・エルミート行列}

\begin{definition}{三つの行列}{tb-symmetric-hermitian}
\[
A^{\mathsf T}=A:\text{実対称},\qquad
A^{\mathsf T}=-A:\text{実交代},\qquad
A^*=A:\text{エルミート}.
\]
\end{definition}

\begin{theorem}{スペクトル定理}{tb-spectral-theorem}
実対称行列は直交行列で、エルミート行列はユニタリ行列で対角化できる。
固有値は全て実数であり、相異なる固有値の固有空間は互いに直交する。
\end{theorem}

ユニタリ対角化では、各固有空間の基底を求め、重複固有値の固有空間内で
Gram--Schmidt 法を使い、得られた正規直交固有ベクトルを列に並べる。

逆に、実行列が直交行列 $Q$ と実対角行列 $D$ により
$A=QDQ^{\mathsf T}$ と書けるなら
$A^{\mathsf T}=QD^{\mathsf T}Q^{\mathsf T}=A$ なので $A$ は対称行列である。

エルミート行列 $A$ で $Av=\lambda v$ とすれば
\[
\lambda\norm v^2=v^*Av=(v^*Av)^*=\overline\lambda\norm v^2
\]
なので $\lambda\in\R$ である。
相異なる固有値 $\lambda,\mu$ の固有ベクトル $v,w$ について
$\lambda\inner{v}{w}=\inner{Av}{w}=\inner{v}{Aw}=\mu\inner{v}{w}$ だから直交する。

非零の実交代行列は実数上で対角化できない。実固有値 $\lambda$ と実固有ベクトル $v$ に対し
\[
\lambda\norm v^2=v^{\mathsf T}Av
=(v^{\mathsf T}Av)^{\mathsf T}=v^{\mathsf T}A^{\mathsf T}v=-v^{\mathsf T}Av
\]
より $\lambda=0$ である。実対角化できるなら全固有値が0なので行列自体が0となり、仮定に反する。

\section{二次形式と行列空間}

\begin{theorem}{Rayleigh 商}{tb-rayleigh}
実対称行列 $A$ の最小・最大固有値を $\lambda_{\min},\lambda_{\max}$ とすると
\[
\lambda_{\min}\le\frac{x^{\mathsf T}Ax}{x^{\mathsf T}x}\le\lambda_{\max}\qquad(x\ne0).
\]
単位球面上の $x^{\mathsf T}Ax$ の最大値・最小値は両端の固有値であり、
等号は対応する固有空間上で成立する。
\end{theorem}

$n$ 次実対称行列全体の次元は、対角成分 $n$ と上三角部分 $n(n-1)/2$ を合わせて
$n(n+1)/2$ である。実交代行列全体の次元は $n(n-1)/2$ である。

\begin{definition}{Frobenius 内積}{tb-frobenius}
同じサイズの実行列に
\[
\inner{A}{B}_{F}=\tr(A^{\mathsf T}B)
\]
と定める。特に列ベクトル $x,y$ に対し
\[
\inner{xx^{\mathsf T}}{yy^{\mathsf T}}_F=(x^{\mathsf T}y)^2.
\]
\end{definition}

$\sum_i a_ix_ix_i^{\mathsf T}=O$ の両辺と $x_jx_j^{\mathsf T}$ の Frobenius 内積を取ると
\[
\sum_i a_i(x_i^{\mathsf T}x_j)^2=0.
\]
係数行列が正則と分かれば全ての $a_i=0$、従って外積行列は一次独立である。
特に $\norm{x_i}=1$, $\abs{\inner{x_i}{x_j}}=r$（$i\ne j$, $0\le r<1$）なら、
$S=\sum_i a_i$ とおくことで各式は
\[
(1-r^2)a_j+r^2S=0
\]
となる。全ての $a_j$ は等しく、元の式から
$[1+(n-1)r^2]a_j=0$、従って全て0である。
3次対称行列空間の次元は6なので、そのような行列は高々6個しか存在できない。
